\begin{table}[htbp]
\caption{Predefined Tests for Attachment-Dependent Policy and Secure-Base Recovery}
\label{tab:h2-h3-tests}
\centering
\scriptsize
\setlength{\tabcolsep}{2pt}
\renewcommand{\arraystretch}{0.85}
\begin{threeparttable}
\begin{tabular}{p{0.09\linewidth}p{0.26\linewidth}p{0.27\linewidth}p{0.26\linewidth}}
\toprule
Test & Model contrast & Target quantity & Decision rule \\
\midrule
H2a & Channel-specific action model without versus with \(\Anx \times e_t\) & Approach/proximity-seeking response to closeness deficit & Retain only if moderation is reproducible and improves held-out prediction \\
H2b & Channel-specific action model without versus with \(\Avo \times e_t\) & Overt approach, distancing, and internal disengagement & Retain channel-specific effects; do not force one signed-action interpretation \\
H2c & Actor-only versus actor--partner moderation & Partner attachment contribution to actor response & Treat as dyadic extension unless incremental fit and prediction justify it \\
H3a & Recovery baseline versus baseline plus \(\Anx,\Avo\) & Recovery explained by insecurity dimensions & Establishes comparison model, not positive security \\
H3b & H3a plus \(\Sec\) & Incremental value of positive secure-base representation & Promotion requires primary-outcome improvement beyond \(\Anx,\Avo\), out-of-sample value, and replication; target retreat alone is not recovery \\
H3c & H3b plus observed/perceived partner support & Whether \(\Sec\) remains informative after proximal support is modeled & Stronger evidence for independent core status; otherwise report conditional attenuation, not identified mediation \\
\bottomrule
\end{tabular}
\begin{tablenotes}[flushleft]
\footnotesize
\item Note. H2 is evaluated on separate action channels as well as any signed summary. H3 is a promotion test for a positive secure-base measure; secondary recovery outcomes cannot rescue failure on the preregistered primary endpoint.
\end{tablenotes}
\end{threeparttable}
\end{table}
